\documentclass{article}\usepackage{booktabs}
\begin{document}
\begin{abstract}
We study sparse routing for efficient inference. Prior gating methods commit to an
expert before the token context resolves, which caps their achievable sparsity. We
defer selection to a single global gate that sees pooled context. On ImageNet this
cuts inference FLOPs by four times at equal accuracy, and an ablation shows the gain
comes from routing rather than added capacity.
\end{abstract}
\section{Introduction}
Efficient inference matters for deployment on edge accelerators and for serving cost.
Given a token stream, previous methods route per layer using a local gate. They work
in many cases. They fail when the discriminative context arrives after the routing
layer, because the gate must commit before that information exists.
We propose a global gate. The basic idea appears in Figure~\ref{fig:teaser}.
Specifically, we pool the sequence, score every expert once, and dispatch. In contrast
to per-layer routing, the decision sees full context. We evaluate on ImageNet and
ADE20K, improving top-1 by 2.4 points at matched FLOPs~\cite{a,b}.
\begin{figure}\includegraphics{t.pdf}\caption{Teaser.}\label{fig:teaser}\end{figure}
\section{Related Work}
Conditional computation routes inputs to subsets of parameters~\cite{c,d,e}. These
assume routing is local. Pruning removes weights statically~\cite{f,g}. Closest to
ours is the layerwise gate of~\cite{h}, which also pools features but commits per
layer.
\section{Method}
We represent the expert set as a bank of MLPs. Given pooled features, we score each
expert and take the top $k$. A remaining problem is load imbalance. We therefore add
an entropy penalty, which keeps utilization flat.
\begin{equation}g = \mathrm{softmax}(Wz)\end{equation}
This yields the dispatch weights used below. The gate costs one extra forward pass and
removes per-layer overhead, reducing routing cost from $O(LN)$ to $O(N)$~\cite{i}.
We train with Adam, learning rate $10^{-3}$, batch size 256, for 300 epochs.
\section{Experiments}
We evaluate on ImageNet-1k and ADE20K against three recent baselines under matched
FLOPs, with three seeds each. Table~\ref{tab:main} reports top-1.
Our method improves top-1 by 2.4 points over the strongest baseline at matched FLOPs.
The gain concentrates on long-tail classes, consistent with rare inputs reaching
specialist experts. On ADE20K the margin narrows to 0.4 points, within seed variance.
Replacing the learned gate with random routing costs 3.1 points, while widening the
MLP to match parameters recovers only 0.4, so the gain comes from routing.
We sweep $k$ from 1 to 8 and find performance flat between 2 and 4.
\begin{table}\caption{Main results.}\label{tab:main}\begin{tabular}{lcc}\toprule
Method & Top-1 & GFLOPs\\\midrule Ours & 85.1 & 4.0\\\bottomrule\end{tabular}\end{table}
\begin{figure}\includegraphics{e.pdf}\caption{Routing entropy per layer.}\label{fig:ent}\end{figure}
See Figure~\ref{fig:ent} for the entropy analysis.
\section{Conclusion}
This paper addressed context-blind routing by deferring expert selection to a global
gate. Experiments show a four times FLOP reduction at equal accuracy. The approach
assumes a fixed expert set; extending to a growing set is the natural next step.
\end{document}
